\section{Reliability y Generalized AI Systems}

Los últimos diez minutos convierten la demo optimista en un problema de sistemas: los agents pueden caer en bucles, son difíciles de probar, pueden quedar sin supervisión y, cuando les falta retroalimentación, producen plan divergence. Después, la clase se apoya en las figuras de LLM OS y de generalized AI system para intentar integrar chat, rules, task engine, router, apps, local index y reflection en una arquitectura unificada.

\subsection{Key Issues: confiabilidad, bucles, evaluación y capacidad de detención}

La slide Key Issues enumera reliability, looping, testing/benchmarking y problemas todavía inmaduros del tipo reset/deployment/observabilidad. La actitud de ingeniería más importante es esta: un agent failure no es una tasa de error única, sino un proceso dinámico en el que una acción modifica las observation y decision posteriores.

\lecturefigure{slide-60-key-agent-issues.jpg}{Problemas clave de confiabilidad de los Autonomous Agents}{Video de la clase de Stanford 00:50:30}

\begin{knowledgebox}{Lectura de la figura: cada uno de los cuatro tipos de problema requiere evidencia distinta}
Reliability requiere la tasa de éxito por tarea y una clasificación de errores; looping requiere step limit y detección de estados repetidos; testing requiere entornos reproducibles y hidden test; oversight requiere trace, kill switch y una persona responsable. Un benchmark que solo evalúa la respuesta final no alcanza a cubrir los sistemas que actúan en entornos reales.
\end{knowledgebox}

\begin{importantbox}{Métricas por niveles de Agent Evaluation}
\begin{tabularx}{\textwidth}{L{0.22\textwidth}X X}
\toprule
Nivel & Métricas de ejemplo & Fallas que no deben ignorarse \\
\midrule
Step & action validity, tool error & hacer clic en el objeto equivocado, parámetros fuera de rango \\
Trajectory & progress, loop rate, budget & ejecución repetida, deriva del estado \\
Task & verified success, human intervention & parece completada, pero la postcondition es incorrecta \\
Safety & unauthorized action, recovery & exceso de privilegios, fuga de información, pérdida irreversible \\
\bottomrule
\end{tabularx}
\end{importantbox}

\teachervoice{El ponente enfatiza que siempre debe ser posible hacer stop al agent; si el sistema sigue en bucle o ejecuta acciones erróneas, las personas necesitan reset, oversight o un kill mechanism. La autonomía no consiste en eliminar la supervisión, sino en trasladarla de la operación paso a paso al control de límites y excepciones.}

\subsection{Plan Divergence: un plan sin retroalimentación se aleja de la realidad}

La slide Plan Divergence representa la desviación con un ideal plan en negro y una agent trajectory en rojo. Al principio el error puede ser pequeño, pero si el sistema sigue planificando únicamente a partir del texto que él mismo generó, sin leer el estado de la página web, los compiler error ni las correcciones del usuario, la desviación se acumula y se convierte en nuevas premisas erróneas. Esta sección vuelve al feedback loop para mostrar cómo usar señales externas verificables para devolver la trajectory al objetivo de la tarea.

\lecturefigure{slide-61-plan-divergence.jpg}{Plan Divergence en ausencia de retroalimentación del entorno}{Video de la clase de Stanford 00:53:09}

\begin{importantbox}{Ver al Agent como un sistema de control por retroalimentación}
Sea $x_t^*$ el estado deseado, $x_t$ el estado real del entorno y $e_t=x_t^*-x_t$ el error. La planificación open-loop, sin retroalimentación, solo ejecuta acciones generadas de antemano; un closed-loop agent lee la observation en cada paso y usa
\[
a_t=\pi(g,x_t,e_t),\qquad x_{t+1}=F(x_t,a_t)
\]
para corregir las acciones posteriores. Lo importante no es que la fórmula sea compleja, sino que $x_t$ provenga de evidencia del entorno y no de lo que el modelo imagina.
\end{importantbox}

\begin{warningbox}{La autorreflexión no sustituye la verificación en el entorno}
Hacer que el mismo modelo lea su propio plan puede corregir contradicciones evidentes, pero no le permite saber si la página web se envió, si el código compila o si el inventario cambió. Conviene priorizar external verifier como compiler, tests, API response, DOM state o receipt; la critique verbal solo sirve como complemento.
\end{warningbox}

\teachervoice{La clase toma el coding agent como ejemplo positivo: el compiler o el IDE devuelven errores explícitos, de modo que el agent sabe «te equivocaste». El ponente atribuye la divergence de los primeros sistemas AutoGPT-style a la falta de feedback confiable; esta explicación está más cerca de la causa raíz que simplemente aumentar las rondas de razonamiento.}

\subsection{LLM OS: construir una capa operativa alrededor del modelo}

La slide LLM OS de Karpathy coloca al LLM en el centro y lo conecta con conversational interface, tools, browser, storage, vision y otras I/O. No es un sistema operativo ya estandarizado, sino una architecture metaphor: el modelo funciona como compute kernel y la capa del sistema se encarga del contexto, los recursos, el enrutamiento y la seguridad.

\lecturefigure{slide-62-llm-os.jpg}{Diagrama de arquitectura de LLM OS de Karpathy presentado en clase}{Video de la clase de Stanford 00:53:33}

\begin{knowledgebox}{Lectura de la figura: qué responsabilidades corresponden a la analogía con un OS}
La chat interface equivale al user shell; context/memory, al working set y al storage; el router, al scheduler; tools/apps, a devices y services; permissions/sandbox, a la protection boundary; reflection/verifier, al error handling. La analogía sirve para detectar componentes faltantes; no implica que el LLM tenga el determinismo ni las garantías de aislamiento de un OS tradicional.
\end{knowledgebox}

\teachervoice{El ponente señala que las personas no siguieron considerando la CPU como «la computadora en sí»; del mismo modo, tampoco debe considerarse siempre al chatbot como la forma final de la IA. El sistema necesita envolver el model con reglas, un motor de tareas, almacenamiento y herramientas para poder asumir trabajos más complejos.}

\subsection{Generalized AI System: chat, rules, routing y reflection}

La slide Generalized AI Systems muestra un flujo de datos más completo: el usuario escribe su entrada en la chat interface; el sistema la combina con las rules para crear una task, que pasa por el task engine y el router para distribuirse a apps/tools o al local index; después, los resultados regresan a reflection y al usuario. Al leer la figura conviene partir del orden de control y no de los nombres de los componentes.

\lecturefigure{slide-63-generalized-ai-systems.jpg}{Boceto de Building Generalized AI Systems presentado en clase}{Video de la clase de Stanford 00:58:51}

\begin{knowledgebox}{Lectura de la figura: una ruta de tarea segura}
La entrada pasa primero por la verificación de reglas y la aclaración del objetivo; el task engine genera un bounded plan y el router elige browser, app, tool o local index; cada ejecutor devuelve evidence y reflection comprueba si se cumple el success predicate. Si no se cumple, el sistema revisa el plan dentro del presupuesto; si alcanza un límite de permisos o de riesgo, devuelve el control al usuario.
\end{knowledgebox}

\begin{warningbox}{Las rules deben restringir el plan antes de la ejecución}
Si las prohibiciones solo se escriben en la plantilla de la respuesta final, el agent puede haber realizado ya una acción sin autorización. Las reglas deben intervenir en plan validation, tool authorization y argument checking; todo plan candidato que viole un principio inmutable debe rechazarse antes de ejecutarse.
\end{warningbox}

\teachervoice{La clase propone las «inherent rules»: si un plan generado por el sistema viola los principios, debe quedar invalidated automáticamente. A continuación, el manager/router decide si accede a browser, apps, tools o local files, y reflection evalúa si la tarea se completó correctamente.}

\subsection{Future Needs: error correction, permissions y sandboxing}

La última slide enumera error correction, security/user permissions, sandboxing y el despliegue de alto riesgo. Con ella, toda la discusión sobre Agents se concreta en los criterios de aceptación más realistas: no basta con completar el camino sin contratiempos; también hay que capturar errores, limitar el acceso, aislar la ejecución y proteger las operaciones business, financial y otras irreversibles.

\lecturefigure{slide-64-future-needs.jpg}{Mecanismos de corrección de errores, seguridad y aislamiento que necesitarán los Agents futuros}{Video de la clase de Stanford 01:00:06}

\begin{importantbox}{Ciclo mínimo de seguridad antes del despliegue}
\begin{enumerate}
  \item \textbf{Least privilege}: cada tool recibe solo los permisos necesarios para completar la tarea actual;
  \item \textbf{Sandbox}: las operaciones de código, archivos y navegación se ejecutan en un entorno restringido;
  \item \textbf{Confirmation gate}: se pide confirmación antes de pagos, eliminaciones, publicaciones, envíos y cambios de cuenta;
  \item \textbf{Audit trail}: se guardan el objetivo, el plan, las acciones, las observaciones, las autorizaciones y los resultados;
  \item \textbf{Recovery}: se admite cancelar, aplicar transacciones compensatorias, revertir versiones y pasar el control a una persona;
  \item \textbf{Evaluation}: se prueban el exceso de privilegios, la inyección y la divergence en escenarios adversarial y en tareas de larga duración.
\end{enumerate}
\end{importantbox}

\teachervoice{Al final, el ponente explica los permissions con un ejemplo muy concreto: el agent puede acceder a DoorDash, pero no necesariamente a la bank account; además, hay que evitar que borre el contenido de la computadora y establecer protecciones ante riesgos comerciales, financieros e irreversibles. No son funciones complementarias para el futuro, sino componentes de un agent que pueda desplegarse.}

\subsection{Resumen de la sección}

La Agent reliability debe evaluarse sobre la trajectory y el environment loop. La plan divergence muestra que la external feedback es necesaria; las figuras de LLM OS y de generalized-system muestran que chat, rules, task engine, router, tools, memory, reflection y permissions deben diseñarse como componentes explícitos; por su parte, la página de future-needs establece error correction, sandboxing e irreversible-action safeguards como requisitos para el despliegue.
